


Sei stets $\mathbb{K}=\mathbb{R}$ oder $\mathbb{K}=\mathbb{C}$.


[[1]]

1. Sei (c) der (ebenfalls durch $\|\cdot\|_{\infty}$ normierte) Unterraum von $\ell^{\infty}=\ell^{\infty}(\mathbb{N})$, welcher aus allen konvergenten Folgen in $\mathbb{K}$ besteht.

[[1.1]]
(i) Zeigen Sie: (c) ist abgeschlossen in $\ell^{\infty}$ und separabel.

[[1.2]]
(ii) Zeigen Sie: Sei $\left(c_0\right)$ der Unterraum von (c) aller Nullfolgen in $\mathbb{K}$, ebenfalls normiert mit $\|\cdot\|_{\infty}$. Geben Sie einen topologischen Isomorphismus $T:(c) \rightarrow\left(c_0\right)$ an.


[[2]]

2. Sei $p \in[1, \infty]$. Für Folgen $a=\left(a_j\right)_j \in \ell^1(\mathbb{Z}) ; x=\left(x_j\right)_j \in \ell^p(\mathbb{Z})$ definieren wir die Faltung
$$
a * x
$$
als eine über $\mathbb{Z}$ indizierte Folge mit den Einträgen
$$
[a * x]_k:=\sum_{j \in \mathbb{Z}} a_{k-j} x_j, \quad k \in \mathbb{Z}
$$


[[2.1]]
(i) Zeigen Sie: Für $a \in \ell^1(\mathbb{Z})$ und $x \in \ell^p(\mathbb{Z})$ ist $a * x \in \ell^p(\mathbb{Z})$ und $\|a * x\|_p \leq\|a\|_1\|x\|_p$.

[[2.2]]
(ii) Zeigen Sie: Die für festes $a \in \ell^1(\mathbb{Z})$ definierte lineare Abbildung
$$
T_a: \ell^p(\mathbb{Z}) \rightarrow \ell^p(\mathbb{Z}), \quad T_a x=a * x
$$
ist stetig mit Operatornorm $\left\|T_a\right\|=\|a\|_1$.


3. 


[[3.1]]
(i) Sei $M$ eine Menge. Zeigen Sie: $\ell^{\infty}(M)$ ist ein Banachraum.


[[3.2]]
(ii) Seien $a, b \in \mathbb{R}, a<b$ und $k \in \mathbb{N}$.

Der $\mathbb{K}$-Vektorraum $C^k[a, b]$ ist definiert als die Menge aller $k$-mal stetig differenzierbaren Funktionen $f:[a, b] \rightarrow \mathbb{K}$. Wir versehen $C^k[a, b]$ mit der Norm $\|\cdot\|_{C^k}$ definiert durch

$$
\|f\|_{C^k}=\sum_{j=0}^k\left\|f^{(j)}\right\|_{\infty} \quad \text { für } f \in C^k[a, b]
$$


Zeigen Sie, dass ( $C^k[a, b],\|\cdot\|_{C^k}$ ) ein Banachraum ist.


[[4.1]]

4. Zeigen Sie durch Angabe eines Gegenbeispiels, dass auf die Voraussetzung $\lim _{n \rightarrow \infty} \operatorname{diam} A_n$ im Satz von Cantor (Satz 3.7. der Vorlesung) nicht verzichtet werden kann.
Können die Mengen $A_n$ in diesem Gegenbeispiel auch als beschränkte Teilmengen eines Banachraums gewählt werden?